\scanpage{133}

et \deleted{\unsure{$R$}} par lemme 1 il suffit de regarder les \uline{extrêmes}. Celui de droite est O.K. \deleted{\unsure{comme}} par \uline{définition}, celui de gauche par hyp de récurrence. \quad O.K.

\uline{Proposition} \uline{4} \quad Soit $i:U\unsure{\text{---}|\longrightarrow}X$ une immersion ouverte, \inserted{avec $X$ excellent, \inserted{de dim $\leq d$},} $F$ un \deleted{\unsure{fais}} $\mathbb Q_\ell$-faisceau \deleted{\unsure{$d$}}$d'$-admissible. Alors les $R^{\unsure{n|-}}i_*(F)$ sont $(d+d'\deleted{\unsure{+2}}\unsure{-1})$-\deleted{\unsure{ad}}admissibles.

\deleted{En effet, soit \inserted{\unsure{de façon qu'on}} $V\underset{j}{\subset}U$, $V$ dense dans $U$, tel que $V$ régulier.} \inserted{Récurrence sur $d+d'$, trivial si $d+d'=0$ car $d=0,d'=0$. Supposons $d+d'>0$, et \deleted{\unsure{soit}} fait si $\delta+\delta'<d+d'$}

On peut supposer $X$ affine donc noeth, et par réc. noeth., \deleted{\unsure{supposons}} le th est démontré pour les $Y\subset X$ fermés, $Y\neq X$, et $X\neq\varnothing$

Soit alors $V\underset{j}{\subset}U$ ouvert dense de $U$, régulier, alors

\diagram{\begin{tikzcd}[ampersand replacement=\&,column sep=small,row sep=small] {} \& Q^*\arrow[dl]\\ F\arrow[rr,no head] \& {} \& \mathbb R j_*F\arrow[ul]\end{tikzcd}}

où $Q$ est concentré sur $Z\cap U$, $Z=\overline{U-V}$ \uline{\unsure{vu|pris}} dans $X$, donc $\neq X$. Donc on est ramené à prouver que les $\mathbb R^n i_*(\mathbb R j_*(F))$ et $\mathbb R^n i_*(Q^*)$ sont admissibles. pour le

% archive sheet number: 132
